% Remap Unicode characters that are missing from Latin Modern fonts.
% Used as --include-in-header by docs/make_pdfs.sh.
% This is a safety net; the primary fix is switching to DejaVu fonts (see
% make_pdfs.sh) which cover most of these natively.
\usepackage{newunicodechar}
\usepackage{amssymb}          % \checkmark, \geq, \approx fallbacks

% pandoc loads unicode-math, which claims these codepoints as math-class
% characters and overrides newunicodechar's text-mode remap at
% \begin{document}. Defining the remaps inside \AtBeginDocument runs them
% after unicode-math finishes, so the substitutions actually take effect
% (otherwise ✓ U+2713 falls through to DejaVu Serif, which lacks the glyph).
\AtBeginDocument{%
  \newunicodechar{≥}{\ensuremath{\geq}}%
  \newunicodechar{≈}{\ensuremath{\approx}}%
  \newunicodechar{γ}{\ensuremath{\gamma}}%
  \newunicodechar{✓}{\ensuremath{\checkmark}}%
  % Emoji the docs use as callout markers. DejaVu has no glyph for either, and
  % a missing glyph is silent in the PDF — the marker simply vanishes, taking
  % the "this is a warning" signal with it. Nearest text-mode equivalents.
  \newunicodechar{⚠}{\textbf{!}}%
  \newunicodechar{💡}{\textbf{*}}%
}

% Long lines in code blocks: wrap them instead of running off the page.
%
% pandoc sets code blocks in fancyvrb's Verbatim, which does not break lines at
% all — a shell command longer than the text width simply runs into the margin
% and off the paper, taking the end of the command with it. Measured on the
% single-file book before this: 155 overfull boxes, the worst 195.6pt (about
% 6.9cm) past the margin.
%
% fvextra extends fancyvrb with line breaking. The Highlighting environment has
% to be redefined rather than configured, because pandoc defines it itself; this
% file is included after pandoc's own definitions, so this one wins. Keep
% commandchars: it is what carries the syntax highlighting.
%
% breakanywhere is on because the offenders here are long paths, URLs and flags
% with no spaces to break at, which breaklines alone cannot help.
%
% breaknonspaceingroup is the one that is easy to miss, and without it the worst
% lines stay broken: pandoc wraps every highlighted run in a macro
% (\StringTok{...}, \NormalTok{...}), and breakanywhere cannot break inside a
% macro argument. A 220-character emcc command in docs/WASM.md still ran 347pt
% past the margin with breaklines+breakanywhere alone; this option is what
% lets a break happen inside the group.
% A fence with no language is not highlighted, and pandoc renders those with
% LaTeX's own verbatim rather than Highlighting — so redefining Highlighting
% alone leaves exactly the blocks that hold long shell commands untouched.
% RecustomVerbatimEnvironment routes plain verbatim through fancyvrb too.
\usepackage{fvextra}
\RecustomVerbatimEnvironment{verbatim}{Verbatim}{breaklines,breakanywhere}
\DefineVerbatimEnvironment{Highlighting}{Verbatim}{%
  breaklines,%
  breakanywhere,%
  breaknonspaceingroup,%
  breaksymbolleft={\raisebox{0.2ex}{\tiny\ensuremath{\hookrightarrow}}},%
  commandchars=\\\{\}%
}

% Long words in ordinary text: break them rather than letting them run off the
% paper. The margin is 25mm (about 71pt), so anything overflowing by more than
% that is not merely ugly — it is off the sheet and unreadable. Measured on the
% book before this: five such boxes, and one of them hid half of the citation
% DOI on page 14.
%
% Two causes, two fixes.
%
% 1. URLs. LaTeX will not break a URL at all by default, so a DOI in running
%    text runs straight off the page. xurl allows a break anywhere in one,
%    without inserting a hyphen. It has to load after hyperref, which the pandoc
%    template loads after this file — hence AtEndPreamble rather than a plain
%    \usepackage here.
\usepackage{etoolbox}
\AtEndPreamble{\usepackage{xurl}}

% 2. Identifiers with underscores, which turn up as example names in narrow
%    table columns (demo_mouse_audio_feedback, Trubutschek_Unconscious_Working_
%    Memory). A hyphen is a legal break point in LaTeX and a literal underscore
%    is not, so a hyphenated name wraps and an underscored one of the same
%    length does not. Allowing a break *after* the underscore adds no character
%    to the text — which matters, because these names get copied.
\AtBeginDocument{%
  \let\goxpyoldunderscore\_%
  \renewcommand{\_}{\goxpyoldunderscore\allowbreak}%
}

% 3. Long inline code — a Go call or a path in a narrow table column. Widening
%    the column cannot fix these: `exp.AddDataVariableNames([]string{...})`
%    measures ~257pt against a 440pt text width, so it would need 58% of the
%    table, and the column beside it holds a 38-character cell with the same
%    problem. The text has to break instead.
%
%    seqsplit permits a break between any two characters, inserting nothing —
%    which is what code needs, since a hyphen would be copied along with the
%    command. It is applied only when the span is too wide for the space it is
%    in (\linewidth is the cell width inside a table), so ordinary short code
%    keeps its normal kerning and is untouched.
%
%    \protected matters: \texttt appears inside section titles, which are
%    written to the .aux as moving arguments. Unprotected, the width test would
%    be expanded there and break the table of contents.
\usepackage{seqsplit}
\newlength{\goxpyttwidth}
\AtBeginDocument{%
  \let\goxpyorigtexttt\texttt
  \protected\def\texttt#1{%
    \settowidth{\goxpyttwidth}{\goxpyorigtexttt{#1}}%
    \ifdim\goxpyttwidth>0.30\linewidth
      \goxpyorigtexttt{\seqsplit{#1}}%
    \else
      \goxpyorigtexttt{#1}%
    \fi
  }%
}
